\section{Increasing Subsequence}
\label{sec:cses-increasing-subsequence}

\problemstatement{Given an array of $n$ integers, find the length of the
longest strictly increasing subsequence (LIS): the longest sequence of
indices $i_1<i_2<\cdots$ with strictly increasing values.}

\begin{patternbox}
\emph{``Longest increasing subsequence''} has an elegant
$O(n\log n)$ solution built on \textbf{patience sorting}: maintain the
smallest possible tail value for an increasing subsequence of each
length, and binary-search where each new element fits.
\end{patternbox}

\subsection*{The tails array}

Keep an array \textit{tails}, where $\textit{tails}[k]$ is the smallest
value that can end an increasing subsequence of length $k+1$ seen so far.
This array is always sorted. For each new value $x$, find the first tail
$\ge x$ (a \texttt{lower\_bound}) and overwrite it with $x$: this either
extends the longest subsequence (if $x$ is larger than every tail,
append) or improves a length-$k$ subsequence's ending to a smaller value,
leaving room for future growth. The LIS length is the final size of
\textit{tails}.

\begin{ideabox}[Keep the Smallest Tail per Length]
Replacing the first tail $\ge x$ with $x$ never shortens any achievable
length and keeps every tail as small as possible, maximising the chance of
future extension. The binary search is what turns the naive $O(n^2)$ DP
into $O(n\log n)$.
\end{ideabox}

\subsection*{Algorithm}

\begin{enumerate}
  \item Start with an empty \textit{tails}.
  \item For each $x$: binary-search the first index with
        $\textit{tails}[i]\ge x$. If none, append $x$ (LIS grew); else
        set $\textit{tails}[i]=x$.
  \item Answer is the size of \textit{tails}.
\end{enumerate}

\subsection*{Dry run}

\dryrun{$a=[7,3,5,3,6,2,9,8]$.}

\begin{itemize}
  \item $7\to[7]$; $3$ replaces $7\to[3]$; $5$ appends $\to[3,5]$; $3$
        replaces $5$? No -- lower\_bound of $3$ is index $0$ ($3\ge3$),
        replace $\to[3,5]$ (unchanged in effect).
  \item $6$ appends $\to[3,5,6]$; $2$ replaces $3\to[2,5,6]$; $9$
        appends $\to[2,5,6,9]$; $8$ replaces $9\to[2,5,6,8]$.
  \item LIS length $=4$ (e.g.\ $3,5,6,9$).
\end{itemize}

\subsection*{C++ implementation}

\begin{lstlisting}
#include <bits/stdc++.h>
using namespace std;

int main() {
    int n;
    cin >> n;
    vector<int> a(n);
    for (int& x : a) cin >> x;

    vector<int> tails;
    for (int x : a) {
        auto it = lower_bound(tails.begin(), tails.end(), x);  // first >= x
        if (it == tails.end()) tails.push_back(x);   // extend LIS
        else *it = x;                                // shrink a tail
    }
    cout << tails.size() << "\n";
    return 0;
}
\end{lstlisting}

\begin{complexitybox}
\textbf{Time Complexity.} $O(n\log n)$ -- one binary search per element.
\textbf{Space Complexity.} $O(n)$ for the tails array. (A simpler
$O(n^2)$ DP with $\textit{dp}[i]=1+\max_{j<i,a_j<a_i}\textit{dp}[j]$ also
solves it for small $n$.)
\end{complexitybox}

\begin{takeawaybox}
LIS in $O(n\log n)$ is the patience-sorting tails trick: keep the minimal
ending value for each length, binary-search each insertion. For strictly
increasing use \texttt{lower\_bound}; for non-decreasing use
\texttt{upper\_bound} -- a one-word change that trips up many.
\end{takeawaybox}
